\begin{lemma}\label{lem:decic77-pure-cusps}
A rational effective divisor of degree $10$ on $X_1(77)$ with
entirely noncuspidal support cannot have entirely cuspidal
reduction at $2$.
\end{lemma}
\begin{proof}
Suppose that $D$ is such a divisor with entirely cuspidal reduction.
Use $\pi:X\to Y$ from Lemma~\ref{lem:decic77-geometry}.
The source cusp signature is $1^{30},3^{10},10^3,30$.
On $Y$, the width-degree pairs are
\[
 (77,1)^6,\qquad(11,3)^2,\qquad(7,2)^3,\qquad(1,6).
\]
A degree-$10$ source cusp divisor therefore pushes forward
either to a divisor on the six rational and two cubic cusps,
or to $5Q$ for one of the three quadratic cusps.
The latter weight is the residue degree ratio $10/2$.
Exact primitive-column orbits verify this complete coverage.
The first kind gives
\[
 \sum_{r=0}^3(r+1)\binom{15-3r}{5}=4989
\]
patterns; the three additional images give $4992$ in total.

For $U=\mathcal C(\mathcal C-1)H^0(Y_{\mathbb F_2},\Omega^1)$,
we have $\dim U=32$.
In one exact differential basis the $69$ restriction blocks
cover rational multiplicities through $10$, cubic multiplicities
through $3$, and multiplicity $5$ at each quadratic cusp.
All $4992$ restrictions have rank $10$.
The rational multiplicity-$10$ and quadratic multiplicity-$5$
blocks are freshly computed; verified deck transformations
transport them to their six and three cusp orbits.
An independent binary replay identifies all $60$ retained
lower-multiplicity kernels across the exact basis transition,
and verifies every restriction rank.
A further native calculation in a fresh model independently
gives rank $10$ at $10P$, for one rational cusp $P$, and at
$5Q$, for one quadratic cusp $Q$.
Deck symmetry covers the six and three corresponding patterns.
Applying Lemma~\ref{lem:decic77-finite-image-Cartier} to
$\pi_*D$ gives the contradiction.
\end{proof}
